\begin{lemma}[Time reversal]
Let $F$ be a finite-dimensional real inner-product space, let $C_0\in\R$ and let
$f$ be a real function on the space-time $\ell^2(F\times\R)$. Let
$\widetilde f$ be its time reversal of \noderef{TimeFlip}, and for a space-time
point $z=(x,t)$ write $\widetilde z=(x,-t)$. Then:
\begin{enumerate}
\item for every $\Omega$, if $f$ is a viscosity subsolution of $F[f]=0$ on
$\Omega$ in the sense of \noderef{IsLimitSubOn}, then $\widetilde f$ is a
viscosity supersolution of $F[\widetilde f]=0$ on
$\widetilde\Omega=\{z:\widetilde z\in\Omega\}$ in the sense of
\noderef{IsLimitSuperOn}, and conversely with sub- and supersolutions
exchanged;
\item for every $\ell$ and every $z$,
\[
  \Penv_\ell\widetilde f(z)=-\Menv_\ell f(\widetilde z),
  \qquad
  \Menv_\ell\widetilde f(z)=-\Penv_\ell f(\widetilde z),
\]
with the envelopes of \noderef{QuadEnvP} and \noderef{QuadEnvM};
\item if $(p,a)$ is a fit for $f$ at $z$ of radius $R$ with constant $C_0$ in the
sense of \noderef{IsFit}, then $(-p,a)$ is a fit for $\widetilde f$ at
$\widetilde z$ of radius $R$ with the same constant; consequently if
$f\in\Aclass_{m,C_0}$ in the sense of \noderef{IsFitClass} then
$\widetilde f\in\Aclass_{m,C_0}$;
\item for every $\kappa$, the function $\widetilde f$ satisfies the horizontal
Jensen bound with constant $\kappa$ on all times, in the sense of
\noderef{HasJensenBound}, if and only if
\[
  f(x,t)-\sum_i\theta_if(y_i,t)\le\kappa r
\]
for every time $t$ and every finite horizontal probability configuration
$(\theta,y)$ with mean $x$ and radius $r$, in the senses of
\noderef{IsBaryConfig} and \noderef{BaryRadius}. Thus time reversal changes the
sign of a horizontal Jensen defect.
\end{enumerate}
\end{lemma}
\begin{proof}
Let $\Pi(x,t)=(x,-t)$, an involutive linear isometry of space-time with
$\Pi^{-1}=\Pi$, so $\widetilde f=-f\circ\Pi$ and
$\widetilde\Omega=\Pi^{-1}(\Omega)$ is open whenever $\Omega$ is.

\emph{(1).} Let $\phi$ be $C^2$ near $z\in\widetilde\Omega$ with
$w\mapsto\widetilde f(w)-\phi(w)$ having a local minimum on $\widetilde\Omega$ at
$z$. Put $\psi=-\phi\circ\Pi$, which is $C^2$ near $\Pi z\in\Omega$. For
$w'\in\Omega$ near $\Pi z$ we have $\Pi w'\in\widetilde\Omega$ near $z$ and
\[
  f(w')-\psi(w')=-\bigl(\widetilde f(\Pi w')-\phi(\Pi w')\bigr)
    \le-\bigl(\widetilde f(z)-\phi(z)\bigr)=f(\Pi z)-\psi(\Pi z),
\]
so $f-\psi$ has a local maximum on $\Omega$ at $\Pi z$, and
\noderef{IsLimitSubOn} gives $F[\psi](\Pi z)\ge0$.

Now compute $F[\psi]$ at $\Pi z$ in terms of $F[\phi]$ at $z$. Since $\Pi$ is
linear with $\Pi^2=\mathrm{Id}$ and $\psi=-\phi\circ\Pi$, the chain rule gives
$\nabla\psi(\Pi z)=-\Pi\,\nabla\phi(z)$ and
$D^2\psi(\Pi z)[\xi]=-D^2\phi(z)[\Pi\xi]$ for every $\xi$. Writing
$\nabla\phi(z)=(\nabla_x\phi,\phi_t)$ we get
$\nabla\psi(\Pi z)=(-\nabla_x\phi,\phi_t)$, so its spatial part is
$-\nabla_x\phi$ and
\[
  F[\psi](\Pi z)=D^2\psi(\Pi z)\bigl[(-\nabla_x\phi,1)\bigr]
   =-D^2\phi(z)\bigl[\Pi(-\nabla_x\phi,1)\bigr]
   =-D^2\phi(z)\bigl[(-\nabla_x\phi,-1)\bigr]
   =-F[\phi](z),
\]
the last step because $A[-\xi]=A[\xi]$ for a bilinear form $A$. Hence
$F[\phi](z)\le0$, which is the supersolution inequality for $\widetilde f$ at
$z$. Continuity of $\widetilde f$ on $\widetilde\Omega$ follows from continuity
of $f$ on $\Omega$. The converse implication is the same computation with maxima
and minima exchanged.

\emph{(2).} Fix $\ell$ and $z=(x,t)$. By \noderef{TimeFlip} and
\noderef{QuadEnvP},
\[
  \Penv_\ell\widetilde f(x,t)
   =\sup_y\left\{-f(y,-t-\ell)-\frac{|y-x|^2}{2\ell}\right\}
   =-\inf_y\left\{f(y,-t-\ell)+\frac{|y-x|^2}{2\ell}\right\}
   =-\Menv_\ell f(x,-t),
\]
using $\sup(-A)=-\inf A$ for a set $A$ of reals, and the time argument
$(-t)-\ell$ in \noderef{QuadEnvM} at the point $\widetilde z=(x,-t)$. The second
identity is the same computation with the roles of $\sup$ and $\inf$ exchanged.

\emph{(3).} Let $(p,a)$ be a fit for $f$ at $z=(x,t)$ of radius $R$. Let $h\in F$
and $\ell\in\R$ with $|\ell|\le R$ and $|h-\ell(-p)|\le R$. Apply
\noderef{IsFit} for $f$ with the increment $(h,-\ell)$: the conditions
$|-\ell|\le R$ and $|h-(-\ell)p|\le R$ hold, so
\[
  \Bigl|f(x+h,t-\ell)-f(x,t)-p\cdot h-\Bigl(a-\frac{|p|^2}2\Bigr)(-\ell)\Bigr|
   \le C_0R .
\]
Multiplying by $-1$ inside the absolute value and using
$\widetilde f(x+h,-t+\ell)=-f(x+h,t-\ell)$, $\widetilde f(\widetilde z)=-f(z)$
and $|-p|=|p|$,
\[
  \Bigl|\widetilde f\bigl(\widetilde z+(h,\ell)\bigr)-\widetilde f(\widetilde z)
    -(-p)\cdot h-\Bigl(a-\frac{|-p|^2}2\Bigr)\ell\Bigr|\le C_0R,
\]
which is \noderef{IsFit} for $\widetilde f$ at $\widetilde z$ of radius $R$ with
parameters $(-p,a)$. If $f\in\Aclass_{m,C_0}$, then by part (1) applied with
$\Omega$ the whole space (whose reflection is the whole space) $\widetilde f$ is
a viscosity solution everywhere, and the fit just constructed exists at every
center and radius because $z\mapsto\widetilde z$ is a bijection; so
$\widetilde f\in\Aclass_{m,C_0}$ by \noderef{IsFitClass}.

\emph{(4).} Fix $\kappa$. By \noderef{TimeFlip}, for a configuration
$(\theta,y)$ with mean $x$ and any time $t$,
\[
  \sum_i\theta_i\widetilde f(y_i,t)-\widetilde f(x,t)
   =f(x,-t)-\sum_i\theta_if(y_i,-t).
\]
Since $t\mapsto-t$ is a bijection of $\R$ and the radius of
\noderef{BaryRadius} does not involve $t$, the condition of
\noderef{HasJensenBound} for $\widetilde f$ with $I=\R$ is equivalent to the
displayed reversed-defect inequality for $f$ at every time and configuration.
\end{proof}
